\section{Introduction}

This chapter defines the formal requirements that guided the design and development of VidChain. The process began with identifying who would actually use the system, what problems they needed solved, and what constraints the hardware environment imposed. Requirements are classified using the MoSCoW method (Must Have, Should Have, Could Have, Won't Have) and are divided into functional and non-functional categories. Hardware constraints are derived directly from the project's technical specifications.

\section{User Personas}

\subsection{The Academic Researcher}

The primary persona for VidChain is an academic researcher operating within high-stakes fields such as computer vision, human-computer interaction, or behavioral forensics. This individual typically manages vast repositories of recorded video—ranging from lab experiment footage and field studies to longitudinal observation archives. Their central objective is to transform these raw pixels into structured, queryable evidence without the prohibitively high temporal cost of manual scrubbing. 

In their current workflow, researchers face a critical "blind spot" where existing tools either demand invasive cloud-based processing (risking sensitive data privacy) or require server-grade infrastructure that is financially and technically inaccessible. Consequently, their analysis often remains shallow, consisting of disjointed observations rather than synthesized narratives. The technical comfort level of this persona is high; they prioritize local-first tools that offer high-precision multimodal extraction—specifically vision, transcriptions, and OCR—within the constraints of standard consumer GPUs (e.g., NVIDIA RTX 3050).

\section{Functional Requirements}

The functional requirements of VidChain define the core capabilities of the IRIS engine, prioritized using the MoSCoW framework. Unlike traditional video players, the system must orchestrate a decentralized pipeline to ensure forensic-grade data grounding.

\begin{table}[H]
\centering
\small
\renewcommand{\arraystretch}{1.3}
\begin{tabularx}{\textwidth}{|l|l|X|}
\hline
\textbf{Ref.} & \textbf{Priority} & \textbf{Description} \\ \hline
FR-1 & \textbf{MUST} & \textbf{Multimodal Ingestion}: Accept local media and process vision, audio (Whisper), and OCR streams simultaneously. \\ \hline
FR-2 & \textbf{MUST} & \textbf{Hybrid Querying}: Enable natural language interaction with the video using combined vector and graph-based retrieval. \\ \hline
FR-3 & \textbf{MUST} & \textbf{Local Execution}: Perform all inference (LLM, Embeddings) locally via OLLAMA/ChromaDB to ensure data privacy. \\ \hline
FR-4 & \textbf{SHOULD} & \textbf{Confidence Scoring}: Self-evaluate generated insights and return a grounding score for each narrative claim. \\ \hline
FR-5 & \textbf{SHOULD} & \textbf{GraphRAG Support}: Build a temporal knowledge graph to resolve multi-hop entity tracking across disjoint scenes. \\ \hline
FR-6 & \textbf{COULD} & \textbf{Semantic Heatmap}: Visualize the temporal distribution of search results across the video timeline for faster navigation. \\ \hline
\end{tabularx}
\caption{MoSCoW Functional Requirements Specification}
\label{tab:moscow_fr}
\end{table}

    \item \textbf{Adaptive Keyframe Filtering.} The pipeline should skip visually redundant frames before they reach heavy compute nodes, reducing GPU load without sacrificing coverage.

    \item \textbf{Hardware Telemetry.} The system should expose live CPU, GPU, and VRAM utilisation during ingestion, visible through the dashboard's Neural HUD.
\end{itemize}

\subsection{Could Have}

\begin{itemize}
    \item \textbf{Emotion Analysis.} The system could include an optional \texttt{EmotionNode} (DeepFace) that classifies the dominant emotion of detected persons on a per-frame basis.

    \item \textbf{Action Classification.} An optional \texttt{ActionNode} (MobileNetV3) could supplement VLM scene descriptions with high-level activity labels (e.g., \textit{suspicious}, \textit{emergency}).

    \item \textbf{Desktop GUI.} A native desktop application (\texttt{vidchain-studio}) could provide an alternative interface for users who prefer a standalone tool over a browser.

    \item \textbf{Report Export.} The system could allow users to export the full conversation history and IRIS responses to a structured Markdown report with a single click.
\end{itemize}

\subsection{Won't Have (in v1.0.0)}

\begin{itemize}
    \item \textbf{Real-time Stream Ingestion.} Live video stream processing (RTSP, webcam) is out of scope for the current release. The pipeline is designed for pre-recorded files only.

    \item \textbf{Multi-user Collaboration.} Shared sessions or concurrent multi-user access are not supported. VidChain is a single-user, local-first tool.

    \item \textbf{Cloud Deployment.} Managed cloud hosting, containerised deployment, or SaaS distribution are explicitly deferred to a future release.
\end{itemize}

\section{Non-Functional Requirements}

Non-functional requirements define the operational constraints and quality attributes that ensure the VidChain engine remains robust under forensic workloads. These are categorized into performance, reliability, and privacy metrics.

\subsection{Performance and Efficiency}
The system is designed for high-efficiency multimodal extraction on consumer hardware. A primary performance requirement is the 60\% reduction in frame processing achieved by the \textbf{AdaptiveKeyframeNode} on static or low-motion footage. During development, this node demonstrated up to 80\% savings, significantly preserving the GPU budget for meaningful visual analysis. Furthermore, the **Moondream VLM** must maintain an inference latency of under 5 seconds per keyframe on the minimum 4\,GB VRAM configuration, a target that was consistently met (2--5 seconds) across the v1.0.0 testing suite. Finally, the **FastAPI gateway** must ensure a query response latency of under 30 seconds for complex RAG operations.

\subsection{Reliability and Fault Tolerance}
Reliability is managed through a decentralized "Fail-Fast but Persist" architecture. Every inference stage, including Whisper, VLM, OCR, and DeepFace, must fail independently; a crash in any single node is prohibited from terminating the entire pipeline. This ensures that the system always returns partial, groundable results to the user. Additionally, the system provides cross-platform media resolution via the **VidChain Media Gateway**, ensuring that video playback functions correctly across both Windows and Unix-based local environments.

\subsection{Privacy and Data Integrity}
VidChain is a local-first system; consequently, no video data, transcripts, or query results are permitted to leave the local network unless the user explicitly configures a remote provider. Security is further enforced through strict session management, where deleting a session triggers an atomic purge of all associated metadata, including ChromaDB vectors, NetworkX knowledge graphs, and JSON archives.

\section{Hardware Constraints}

The system's architectural trade-offs are driven by the requirement to run on consumer hardware, particularly laptops with limited VRAM.

\begin{table}[H]
\centering
\small
\caption{Hardware Requirements for VidChain v1.0.0}
\label{tab:hardware}
\begin{tabular}{|p{2.5cm}|>{\raggedright\arraybackslash}p{4.5cm}|>{\raggedright\arraybackslash}p{4.5cm}|}
\hline
\rowcolor[HTML]{EFEFEF}
\textbf{Component} & \textbf{Minimum Configuration} & \textbf{Optimised Configuration} \\ \hline
\textbf{RAM} & 16\,GB & 32\,GB \\ \hline
\textbf{GPU VRAM} & 4\,GB (NVIDIA RTX 3050+) & 12\,GB (NVIDIA RTX 4070+) \\ \hline
\textbf{CUDA Version} & 12.1+ & 12.1+ \\ \hline
\textbf{Python} & 3.11+ & 3.11 or 3.12 \\ \hline
\textbf{Storage} & \textasciitilde{}10\,GB (models + DB) & 20\,GB+ (multiple sessions) \\ \hline
\end{tabular}
\end{table}

\subsection{VRAM Budgeting and Optimization}
The 4\,GB VRAM target informed several critical design decisions. To honor this constraint, VidChain utilizes \textbf{FP32 precision} for the VisionEngine to avoid the FP16 instabilities common in 30-series laptop GPUs, while simultaneously offloading the \textbf{EmotionNode (DeepFace)} to the CPU to prevent VRAM competition. The choice of \textbf{Moondream} as the default VLM was pivotal, as it delivers high-fidelity scene understanding at 186$\times$ lower latency than standard 7B-parameter models on edge hardware. These optimizations, combined with the **Adaptive Keyframe filtering**, ensure that the IRIS engine remains operational even during peak inference loads on consumer-grade machines.

\section{Summary}

This chapter established the requirements boundary for VidChain v1.0.0. The single primary persona, the academic researcher, shaped the emphasis on local-first privacy, consumer hardware compatibility, and natural language accessibility. The MoSCoW classification separates the core intelligence pipeline (Must Have) from optional enrichment nodes (Could Have), providing a clear scope boundary. Non-functional requirements formalise the performance targets and reliability guarantees that were validated iteratively during development, and hardware constraints reflect real architectural trade-offs made to keep the system accessible on edge hardware.

\section{System Models}

Two diagrams are provided to complement the written requirements: a Use Case Diagram capturing the interactions between the primary persona and the system, and a Level-1 Data Flow Diagram tracing how data moves through each node in the pipeline.

\subsection{Use Case Diagram}

Figure~\ref{fig:use_case} illustrates the full set of interactions available to the Academic Researcher persona. The core use cases, ingesting a video, querying content, and viewing a summary, each include a subordinate system process: the node pipeline, GraphRAG retrieval, and the Map-Reduce engine respectively. This makes explicit that the user-facing actions are backed by non-trivial internal workflows, not simple database lookups.

\begin{figure}[H]
\centering
\includegraphics[width=0.90\textwidth]{Use_Case.pdf}
\caption{Use Case Diagram --- Functional Requirements and Actor Interactions}
\label{fig:use_case}
\end{figure}